\section{Introduction}
\label{sec:intro}

Understanding when and why neural networks generalize — particularly in
regimes where training and generalization dynamics decouple — remains one
of the most consequential open problems in deep learning theory. A
striking instance of this decoupling is the \emph{grokking} phenomenon,
first described by Power et al.~\cite{power2022grokking}: a small
transformer trained on modular arithmetic tasks overfits the training
data for thousands of epochs before suddenly achieving near-perfect
generalization on held-out examples, long after any conventional
early-stopping criterion would have terminated training. Existing
explanations approach this transition from complementary angles. Nanda
et al.~\cite{nanda2023progress} identify discrete Fourier transform
circuits that form gradually and eventually dominate the network's
behavior, framing grokking as a smooth internal competition between
memorizing and generalizing subnetworks. Hwang and
Park~\cite{hwang2026intrinsic} propose that intrinsic task symmetries
drive a three-stage representational reorganization culminating in
task-aligned geometry. Both perspectives operate at the level of
specific circuits or algebraic symmetry groups, leaving open the
question of whether higher-level topological properties of the full
representation space follow a consistent, measurable pattern across
grokking events.

In this work, we address this question by applying Topological Data
Analysis (TDA) to track global topological invariants — persistent
homology (Betti numbers) and intrinsic dimensionality — across the
layer-wise activation spaces of a small transformer throughout training.
TDA is well-suited to this setting: it operates directly on the shape
of high-dimensional point clouds without assuming a parametric structure,
and its core summary statistics are provably stable under small
perturbations of the input data~\cite{cohen2005stability}. Our main
contributions are: (i) a modular TDA pipeline combining Maximum
Likelihood intrinsic dimension estimation~\cite{levina2004maximum},
topology-guided dimensionality reduction, and persistent homology via
GUDHI~\cite{gudhi2025}; (ii) a layer-wise topological characterization
of grokking showing that $\beta_0$ decreases and stabilizes in decoder
layers while converging to a higher stable value in the linear layer,
concurrently with the rise in validation accuracy; (iii) evidence that
intrinsic dimension decreases consistently across all layers in
generalizing conditions and remains elevated in non-generalizing ones;
and (iv) a comparative analysis across five training fractions and two
modular arithmetic tasks that identifies topological markers
distinguishing generalization from persistent memorization. Our analysis
is scoped to a single small transformer architecture trained on modular
arithmetic; the extent to which these findings generalize to larger
models or natural language tasks remains an open question. The identified
relationships between topological structure and generalization are
correlational, not causal.